\chapter{Testes e Avaliação}
\label{chap:testes}

\section{Introdução}
\label{sec:test-intro}

Este capítulo descreve os testes realizados para avaliar o sistema Gold Lock em três dimensões: usabilidade (testes com utilizadores reais), performance técnica e precisão do módulo de \ac{ML}. A avaliação segue as métricas definidas nos requisitos não-funcionais (Secção~\ref{sec:req-nao-funcionais}).

% TODO: Preencher com resultados reais após os testes (Sprints 9-10)


\section{Testes de Usabilidade}
\label{sec:test-usabilidade}

\subsection{Metodologia}

Os testes de usabilidade foram conduzidos com um grupo de 10 a 15 participantes, recrutados entre colegas da \ac{UBI}, familiares e amigos. Cada sessão de teste incluiu:

\begin{enumerate}
    \item Breve introdução ao sistema (2 minutos)
    \item Execução de 5 a 7 tarefas predefinidas
    \item Preenchimento do questionário \ac{SUS} (System Usability Scale)
    \item Entrevista breve sobre a experiência (5 minutos)
\end{enumerate}

\subsection{Tarefas de Teste}

As tarefas definidas para os testes de usabilidade foram:

\begin{enumerate}
    \item Registar uma conta e completar o onboarding
    \item Ligar uma conta bancária (ambiente sandbox)
    \item Consultar o saldo total e os gastos por categoria no dashboard
    \item Corrigir a categoria de uma transação
    \item Criar um orçamento mensal para a categoria ``Restauração''
    \item Definir uma meta de poupança
    \item Executar uma simulação de \ac{IRS}
\end{enumerate}

\subsection{Resultados SUS}

% TODO: Preencher com resultados reais
\begin{table}[H]
\centering
\caption{Resultados do questionário SUS (a preencher).}
\label{tab:sus-results}
\begin{tabularx}{\textwidth}{l c c c}
\toprule
\textbf{Métrica} & \textbf{Objetivo} & \textbf{Resultado} & \textbf{Avaliação} \\
\midrule
Pontuação SUS média & $> 70$ & --- & --- \\
Tempo médio por tarefa & $< 30$s & --- & --- \\
Taxa de sucesso de tarefas & $> 85\%$ & --- & --- \\
\bottomrule
\end{tabularx}
\end{table}


\section{Testes de Performance}
\label{sec:test-performance}

Os testes de performance foram realizados utilizando Lighthouse (Google Chrome) para métricas de frontend e ferramentas de load testing para o backend.

% TODO: Preencher com resultados reais
\begin{table}[H]
\centering
\caption{Resultados dos testes de performance (a preencher).}
\label{tab:perf-results}
\begin{tabularx}{\textwidth}{l c c c}
\toprule
\textbf{Métrica} & \textbf{Objetivo} & \textbf{Resultado} & \textbf{Avaliação} \\
\midrule
Tempo de carregamento inicial & $< 2$s & --- & --- \\
Tempo de resposta \ac{API} (p95) & $< 500$ms & --- & --- \\
FPS nas animações & $> 55$ FPS & --- & --- \\
Lighthouse Performance Score & $> 90$ & --- & --- \\
\bottomrule
\end{tabularx}
\end{table}


\section{Avaliação do Modelo de ML}
\label{sec:test-ml}

A avaliação do modelo de categorização automática de transações foi realizada com um dataset de teste separado do dataset de treino.

% TODO: Preencher com resultados reais
\begin{table}[H]
\centering
\caption{Métricas de avaliação do modelo de ML (a preencher).}
\label{tab:ml-results}
\begin{tabularx}{\textwidth}{l c c c}
\toprule
\textbf{Métrica} & \textbf{Objetivo} & \textbf{Resultado} & \textbf{Avaliação} \\
\midrule
Accuracy & $> 80\%$ & --- & --- \\
Precision (macro) & $> 75\%$ & --- & --- \\
Recall (macro) & $> 75\%$ & --- & --- \\
F1-Score (macro) & $> 75\%$ & --- & --- \\
Tempo por transação & $< 200$ms & --- & --- \\
\bottomrule
\end{tabularx}
\end{table}


\section{Validação do Simulador IRS}
\label{sec:test-irs}

O simulador de \ac{IRS} foi validado contra 10 cenários fiscais reais com resultados conhecidos, cobrindo diferentes escalões de rendimento, estados civis e tipos de deduções.

% TODO: Tabela com cenários de teste e resultados


\section{Conclusões}
\label{sec:test-conclusoes}

% TODO: Preencher após a realização dos testes
Os resultados dos testes serão analisados nesta secção, identificando os pontos fortes e as áreas de melhoria do sistema.
